\scanpage{20}
lorsque $X$ et $Y$ sont des courbes \uline{algébriques}
anabéliennes sur $\mathbb C$.

Pourtant, il \deleted{doit être} \inserted{est} possible tout au
moins, sur un corps alg. \inserted{clos $k$} quelconque,
de définir un \uline{foncteur} qui va du
\uline{groupoïde} de Teichmüller des
courbes de type $(g,\nu)$, \inserted{(disons)} avec un
paquet \deleted{de} $I$ de cardinal $\nu'\leq\nu$ de
pts épinglés par $I$, vers le \uline{groupoïde}
de Teichmüller de type $(g,\nu-\nu')$,
par ``bouchage des trous en $I$'', \deleted{\unsure{cela} dit, si}
avec \unsure{une|les} \unsure{propriété|propriétés} de transitivi-
té \unsure{évidentes|évidents}, pour des bouchages de
trous successifs. Mais je vois malaisé-
ment comment l'inclure dans une
notion correspondante de ``morphismes d'
inclusion'', qui pourrait se composer
avec les morphismes dont on dispose
déjà dans la \uline{catégorie} (pas groupoïde !)
de Galois-Teichmüller. À moins d'y
\margin{19}
